% OpenStax Calculus Volume 2 -- Section 5.4 Exercises: Comparison Tests
% Exercises reproduced from OpenStax, licensed under CC BY 4.0.
% Source: https://openstax.org/books/calculus-volume-2/pages/5-4-comparison-tests
\documentclass{exercises}
\title{OpenStax Calculus Volume 2}
\subtitle{Section 5.4 Exercises: Comparison Tests}
\sourceurl{https://openstax.org/books/calculus-volume-2/pages/5-4-comparison-tests}
\author{}
\date{}

\begin{document}
\maketitle


Use the comparison test to determine whether the following series converge.

\begin{enumerate}[start=1]
\item $\sum_{n=1}^{\infty}a_{n}$ where $a_{n}=\frac{2}{n(n+1)}$
\item $\sum_{n=1}^{\infty}a_{n}$ where $a_{n}=\frac{1}{n(n+1/2)}$
\item $\sum_{n=1}^{\infty}\frac{1}{2(n+1)}$
\item $\sum_{n=2}^{\infty}\frac{1}{2(n-1)}$
\item $\sum_{n=2}^{\infty}\frac{1}{{(n\ln n)}^{2}}$
\item $\sum_{n=1}^{\infty}\frac{n!}{(n+2)!}$
\item $\sum_{n=1}^{\infty}\frac{1}{n!}$
\item $\sum_{n=1}^{\infty}\frac{\sin(1/n)}{n^{2}}$
\item $\sum_{n=1}^{\infty}\frac{\sin^{2}n}{n^{2}}$
\item $\sum_{n=1}^{\infty}\frac{\sin(1/n)}{(\sqrt{n})^{3}}$
\item $\sum_{n=1}^{\infty}\frac{n^{1.2}-1}{n^{2.3}+1}$
\item $\sum_{n=1}^{\infty}\frac{\sqrt{n+1}-\sqrt{n}}{n}$
\item $\sum_{n=1}^{\infty}\frac{\sqrt[4]{n}}{\sqrt[3]{n^{4}+n^{2}}}$
\end{enumerate}

Use the limit comparison test to determine whether each of the following series converges or diverges.

\begin{enumerate}[resume]
\item $\sum_{n=1}^{\infty}\left(\frac{\ln n}{n}\right)^{2}$
\item $\sum_{n=1}^{\infty}\left(\frac{\ln n}{n^{0.6}}\right)^{2}$
\item $\sum_{n=1}^{\infty}\frac{\ln(1+\frac{1}{n})}{n}$
\item $\sum_{n=1}^{\infty}\ln(1+\frac{1}{n^{2}})$
\item $\sum_{n=1}^{\infty}\frac{1}{4^{n}-3^{n}}$
\item $\sum_{n=1}^{\infty}\frac{1}{n^{2}-n\sin n}$
\item $\sum_{n=1}^{\infty}\frac{1}{e^{(1.1)n}-3^{n}}$
\item $\sum_{n=1}^{\infty}\frac{1}{e^{(1.01)n}-3^{n}}$
\item $\sum_{n=1}^{\infty}\frac{1}{n^{1+1/n}}$
\item $\sum_{n=1}^{\infty}\frac{1}{2^{1+1/n}n^{1+1/n}}$
\item $\sum_{n=1}^{\infty}\left(\frac{1}{n}-\sin\left(\frac{1}{n}\right)\right)$
\item $\sum_{n=1}^{\infty}\left(1-\cos\left(\frac{1}{n}\right)\right)$
\item $\sum_{n=1}^{\infty}\frac{1}{n}\left(\frac{\pi}{2}-\arctan n\right)$
\item $\sum_{n=1}^{\infty}\left(1-\frac{1}{n}\right)^{n}$ (Hint: $\left(1-\frac{1}{n}\right)^{n}\to 1/e$.)
\item $\sum_{n=1}^{\infty}(1-e^{-1/n})$ (Hint: $1/e\approx(1-1/n)^{n}$, so $1-e^{-1/n}\approx 1/n$.)
\item Does $\sum_{n=2}^{\infty}\frac{1}{{(\ln n)}^{p}}$ converge if $p$ is large enough? If so, for which $p$?
\item Does $\sum_{n=1}^{\infty}{(\frac{(\ln n)}{n})}^{p}$ converge if $p$ is large enough? If so, for which $p$?
\item For which $p$ does the series $\sum_{n=1}^{\infty}2^{pn}/3^{n}$ converge?
\item For which $p>0$ does the series $\sum_{n=1}^{\infty}\frac{n^{p}}{2^{n}}$ converge?
\item For which $r>0$ does the series $\sum_{n=1}^{\infty}\frac{r^{n^{2}}}{2^{n}}$ converge?
\item For which $r>0$ does the series $\sum_{n=1}^{\infty}\frac{2^{n}}{r^{n^{2}}}$ converge?
\item Find all values of $p$ and $q$ such that $\sum_{n=1}^{\infty}\frac{n^{p}}{{(n!)}^{q}}$ converges.
\item Does $\sum_{n=1}^{\infty}\frac{n\pi}{2}$ converge or diverge? Explain.
\item Explain why, for each $n$, at least one of $\{|\sin n|,|\sin(n+1)|,\ldots,|\sin(n+6)|\}$ is larger than $1/2$. Use this relation to test convergence of $\sum_{n=1}^{\infty}\frac{|\sin n|}{\sqrt{n}}$.
\item Suppose that $a_{n}\ge 0$ and $b_{n}\ge 0$ and that $\sum_{n=1}^{\infty}a_{n}^{2}$ and $\sum_{n=1}^{\infty}b_{n}^{2}$ converge. Prove that $\sum_{n=1}^{\infty}a_{n}b_{n}$ converges and $\sum_{n=1}^{\infty}a_{n}b_{n}\le \frac{1}{2}\left(\sum_{n=1}^{\infty}a_{n}^{2}+\sum_{n=1}^{\infty}b_{n}^{2}\right)$.
\item Does $\sum_{n=2}^{\infty}2^{-\ln \ln n}$ converge? (Hint: Write $2^{\ln \ln n}$ as a power of $\ln n$.)
\item Does $\sum_{n=2}^{\infty}(\ln n)^{-\ln n}$ converge? (Hint: Use $n=e^{\ln n}$ to compare to a $p$-series.)
\item Does $\sum_{n=2}^{\infty}(\ln n)^{-\ln \ln n}$ converge? (Hint: Compare $a_{n}$ to $1/n$.)
\item Show that if $a_{n}\ge 0$ and $\sum_{n=1}^{\infty}a_{n}$ converges, then $\sum_{n=1}^{\infty}a_{n}^{2}$ converges. If $\sum_{n=1}^{\infty}a_{n}^{2}$ converges, does $\sum_{n=1}^{\infty}a_{n}$ necessarily converge?
\item Suppose that $a_{n}>0$ for all $n$ and that $\sum_{n=1}^{\infty}a_{n}$ converges. Suppose that $b_{n}$ is an arbitrary sequence of zeros and ones. Does $\sum_{n=1}^{\infty}a_{n}b_{n}$ necessarily converge?
\item Suppose that $a_{n}>0$ for all $n$ and that $\sum_{n=1}^{\infty}a_{n}$ diverges. Suppose that $b_{n}$ is an arbitrary sequence of zeros and ones with infinitely many terms equal to one. Does $\sum_{n=1}^{\infty}a_{n}b_{n}$ necessarily diverge?
\item Complete the details of the following argument: If $\sum_{n=1}^{\infty}\frac{1}{n}$ converges to a finite sum $s$, then $\frac{1}{2}s=\frac{1}{2}+\frac{1}{4}+\frac{1}{6}+\cdots$ and $s-\frac{1}{2}s=1+\frac{1}{3}+\frac{1}{5}+\cdots$. Why does this lead to a contradiction?
\item Show that if $a_{n}\ge 0$ and $\sum_{n=1}^{\infty}a_{n}^{2}$ converges, then $\sum_{n=1}^{\infty}\sin^{2}(a_{n})$ converges.
\item Suppose that $a_{n}/b_{n}\to 0$ in the comparison test, where $a_{n}\ge 0$ and $b_{n}\ge 0$. Prove that if $\sum b_{n}$ converges, then $\sum a_{n}$ converges.
\item Let $b_{n}$ be an infinite sequence of zeros and ones. What is the largest possible value of $x=\sum_{n=1}^{\infty}\frac{b_{n}}{2^{n}}$?
\item Let $d_{n}$ be an infinite sequence of digits, meaning $d_{n}$ takes values in $\{0,1,\ldots,9\}$. What is the largest possible value of $x=\sum_{n=1}^{\infty}\frac{d_{n}}{10^{n}}$ that converges?
\item Explain why, if $x>1/2$, then $x$ cannot be written as $x=\sum_{n=2}^{\infty}\frac{b_{n}}{2^{n}}$, where $b_{n}=0$ or $1$ (and $b_{1}=0$).
\item \techrequired Evelyn has a perfect balancing scale, an unlimited number of $1\,\text{kg}$ weights, and one each of $1/2\,\text{kg}$, $1/4\,\text{kg}$, $1/8\,\text{kg}$, and so on. She wishes to weigh a meteorite of unspecified origin to arbitrary precision. Assuming the scale is big enough, can she do it? What does this have to do with infinite series?
\item \techrequired Robert wants to know his body mass to arbitrary precision. He has a big balancing scale that works perfectly, an unlimited collection of $1\,\text{kg}$ weights, and nine each of $0.1\,\text{kg}$, $0.01\,\text{kg}$, $0.001\,\text{kg}$, and so on. Assuming the scale is big enough, can he do this? What does this have to do with infinite series?
\item The series $\sum_{n=1}^{\infty}\frac{1}{2n}$ is half the harmonic series and hence diverges. It is obtained from the harmonic series by deleting all terms in which $n$ is odd. Let $m>1$ be fixed. Show, more generally, that deleting all terms $1/n$ where $n=mk$ for some integer $k$ also results in a divergent series.
\item In view of the previous exercise, it may be surprising that a subseries of the harmonic series in which about one in every five terms is deleted might converge. A depleted harmonic series is a series obtained from $\sum_{n=1}^{\infty}\frac{1}{n}$ by removing any term $1/n$ if a given digit, say $9$, appears in the decimal expansion of $n$. Argue that this depleted harmonic series converges by answering the following questions.
\begin{enumerate}[label=(\alph*)]
  \item How many whole numbers $n$ have $d$ digits?
  \item How many $d$-digit whole numbers, denoted by $h(d)$, do not contain $9$ as one or more of their digits?
  \item What is the smallest $d$-digit number $m(d)$?
  \item Explain why the depleted harmonic series is bounded by $\sum_{d=1}^{\infty}\frac{h(d)}{m(d)}$.
  \item Show that $\sum_{d=1}^{\infty}\frac{h(d)}{m(d)}$ converges.
\end{enumerate}
\item Suppose that a sequence of numbers $a_{n}>0$ has the property that $a_{1}=1$ and $a_{n+1}=\frac{1}{n+1}S_{n}$, where $S_{n}=a_{1}+\cdots+a_{n}$. Can you determine whether $\sum_{n=1}^{\infty}a_{n}$ converges? (Hint: $S_{n}$ is monotone.)
\item Suppose that a sequence of numbers $a_{n}>0$ has the property that $a_{1}=1$ and $a_{n+1}=\frac{1}{(n+1)^{2}}S_{n}$, where $S_{n}=a_{1}+\cdots+a_{n}$. Can you determine whether $\sum_{n=1}^{\infty}a_{n}$ converges? (Hint: $S_{2}=a_{2}+a_{1}=a_{2}+S_{1}=a_{2}+1=1+1/4=(1+1/4)S_{1}$, $S_{3}=\frac{1}{3^{2}}S_{2}+S_{2}=(1+1/9)S_{2}=(1+1/9)(1+1/4)S_{1}$, etc. Look at $\ln(S_{n})$, and use $\ln(1+t)\le t$ for $t>0$.)
\end{enumerate}

\end{document}
